\section{Lean formalization: checked core and a complete path forward}
\label{sec:formalization}

\subsection{Pinned object syntax and standalone build}

The archive builds with Lean 4.32.0 and its standard library. It includes
unchanged copies of the pinned repository's \code{Syntax.lean},
\code{Deduction.lean}, \code{PR.lean} and
\code{NativeBiInterpretation.lean}. The first three module bodies are
concatenated into \code{ObjectKernel.lean}; their imports are replaced
by \code{import Std}, and a compatibility notation spells the
natural-number type. No object definition, axiom constructor, inference
rule or upstream theorem body is changed. A reproducible provenance
checker verifies the embedded bodies and the pinned Git blob hashes.
The full native bi-interpretation file is retained as a reference and
is not rebuilt by the standalone core.

This arrangement avoids downloading Mathlib merely to use the small
deductive kernel. Integration into ProveIt replaces the compatibility
import by the repository's \code{ConciseIPRA.PR} import and retains
the new namespaces and proofs. The validation record is a standalone
build against the verified upstream bodies, not a claim that the whole
ProveIt workspace has been rebuilt.

\begin{longtable}{P{.29\linewidth}P{.40\linewidth}P{.21\linewidth}}
\toprule
Module & Principal content & Status\\
\midrule\endhead
\path{ObjectKernel.lean} & The three upstream syntax, deduction and
arithmetic module bodies. & Compiled; provenance checked.\\[3pt]
\path{ObjectFusion.lean} & PR iteration equations and the fusion
derivation builder. & Compiled object theorems.\\[3pt]
\path{ObjectSimulation.lean} & Substitution preserves QF status;
relational simulation builder. & Compiled object theorem.\\[3pt]
\path{SquaringTransport.lean} & Repeated-squaring transport and an
unconditional successor instance. & Compiled object theorems.\\[3pt]
\path{PrimitiveEigenvalue.lean} & Finite weighted eigen-equations
and the primitive-vector denominator lemma. & Compiled host integer
theorems.\\[3pt]
\path{EigenvalueBound.lean} & Finite absolute-column estimate and
the eigenvalue bound. & Compiled host integer theorems.\\[3pt]
\path{FiniteAlgebraTargets.lean} & Complete finite programs,
remaining theorem specifications and the filter reduction. & Definitions
compile; main completeness target open.\\[3pt]
\path{FiniteAlgebraExamples.lean} & Exact execution of the elliptic
section certificate and grid. & Compiled diagnostics.\\
\bottomrule
\end{longtable}

There are no custom axioms, opaque placeholders, \code{sorry} or
\code{admit} declarations in the new modules. The target file does
not smuggle the desired theorem into a structure field: the remaining
claims are explicitly named propositions and are not asserted.

\subsection{The completed syntactic theorems}

The important distinction is the type of a theorem's conclusion. In
namespace \path{LeanProofs.ConciseIPRA.FiniteMethods}, the simplified
shape of the checked fusion theorem is
\begin{verbatim}
theorem fusion (f g h : PR 1)
    (hlocal : ([] : List (Fm PTm 1)) |-[PRAAx]
      commutes f g h) :
    ([] : List (Fm PTm 1)) |-[PRAAx]
      Fm.all (fusionInvariant f g h)
\end{verbatim}
Here the printed ASCII turnstile is explanatory; the source uses the
repository's scoped judgment notation. \code{commutes} is precisely
$h(f(x))=g(h(x))$, and \code{fusionInvariant} is the equality in
Theorem~\ref{thm:fusion}. No extra assumption about an evaluation
function is present.

For simulation the inputs are $R:\code{Fm PTm 2}$, a proof of
\code{R.QF}, $f,g:\code{PR 1}$, and the local judgment
$[R]\vdash_{\code{PRAAx}}R(f(x),g(y))$. The output is
$[R]\vdash_{\code{PRAAx}}\forall n\,R(\iter_f(n,x),\iter_g(n,y))$.
The code verifies every lifting of a substitution under a binder.
Both modules use \code{PRA.qfInd}, \code{Prf.subst},
\code{Prf.eqSubst} and the defining equations, exactly as the human
proofs specify.

The type \code{Prf} lives in Lean's proposition universe; this
development proves derivability and supplies a checked derivation term.
It does not serialize a numerical proof trace for the tiny Python
checkers in the other project. Such serialization, if desired, is
another implementation task. This distinction has no effect on the
object judgment proved here.

\subsection{The integer arithmetic already checked}

\code{PrimitiveEigenvalue.lean} stores a finite list of triples
$(e_i,w_i,z_i)$, representing a Bezout coefficient, a primitive-vector
coordinate, and its integral matrix image. It defines
\[
          W=\sum_i e_iw_i,\qquad Z=\sum_i e_iz_i.
\]
From the finite eigen-equations $bz_i=aw_i$ it proves $bZ=aW$.
If $W=1$, then $b\mid a$. A supplied reducedness certificate
$ua+vb=1$ and $b>0$ force $b=1$. The proof is finite integer algebra,
not a call to a spectral theorem.

Equivalently, after primitive scaling the human argument may simply
multiply $z_i=\lambda w_i$ by $e_i$ and add to get
$\lambda=\sum_i e_i z_i\in\Z$. The denominator formulation
matches the explicit integer-coded implementation conveniently.

\code{EigenvalueBound.lean} defines the dot product with one matrix
column and its absolute sum by list recursion. It proves
\[
 \left|\sum_iw_i a_i\right|\le W\sum_i|a_i|
       \quad\hbox{if every }|w_i|\le W.
\]
Given a selected coordinate of absolute value $W>0$ and its eigenvalue
equation, it proves $|\lambda|\le R$ when that column sum is at
most $R$. Constructing the selected coordinate by a maximum scan,
and constructing the primitive vector from rational evaluations, are
explicit remaining host lemmas; neither is hidden in the current
theorem statement.

\subsection{Complete definitions for the rational target}

The target module uses \code{Monomial n} with a rational coefficient
and a length-$n$ natural exponent vector, and
\code{Polynomial n := List (Monomial n)}. Repeated monomials are
allowed. \code{coefficientAt} adds all coefficients with a given
exponent vector. \code{polynomialEq} compares these sums on the
union of the two finite supports; outside that union both sums are zero.
Polynomial addition, subtraction, multiplication and evaluation are
fully defined finite list programs.

A \code{Certificate n s d} contains the $d$ generators, unit
coefficients, unit ideal witnesses, $n$ positive integer denominators,
$n$ integral action matrices and their ideal witnesses. The Lean
representation stores the already-cleared identities
\[
 \Delta_kX_kb_j-\sum_i(A_k)_{ij}b_i
                     =\sum_\ell H_{kj\ell}F_\ell.
\]
The Python representation initially stores rational $M_k$ and then
computes $\Delta_k$; to convert it, multiply its residual witnesses
by $\Delta_k$. Matrix storage order differs explicitly: Python stores
rows; the Lean target stores columns. The mathematical convention is
the same row-vector action $vM$.

The executable target definitions then compute absolute column sums,
the rational coordinate grids, their Cartesian product, the common-zero
test, and its finite filter. The main target is
\begin{verbatim}
def EnumerationGoal : Prop :=
  forall n s d (F : Vector (Polynomial n) s)
      (C : Certificate n s d) (p : Vector Rat n),
    certificateOK F C = true ->
      (p.toList in enumerate F C <->
       zeroTest F p.toList = true)
\end{verbatim}
Again the listing is a legible ASCII rendering; the source is valid
Lean. Arity is enforced by vectors, and the executable zero test
also checks the length of its list input. Invalid certificates produce
an empty result in the target function; completeness is asserted only
when the checker accepts.

\subsection{A noncircular completion plan}

The remaining host proof has the following concrete dependency order.
\begin{enumerate}
\item \textbf{Finite polynomial semantics.} Prove evaluation preserves
addition, scaling and multiplication. Prove that collecting duplicate
monomials preserves evaluation. These finite distributive identities
establish \code{EvaluationTransportGoal}.
\item \textbf{Certificate substitution.} Evaluate each accepted
coefficient identity at the fixed common zero. The ideal-combination
side vanishes by the previous multiplicativity lemma. This proves
\code{CertificateSoundnessGoal} and extracts the unit and action laws.
\item \textbf{Primitive scaling.} Compute a common denominator of the
nonzero evaluation row; compute its gcd and Bezout coefficients. Prove
the resulting integer row is primitive and preserves the eigen-equations.
Instantiate the checked denominator lemma for each coordinate.
\item \textbf{The maximum scan.} Track the first maximizing index of a
nonzero integer row. Prove its entry has positive absolute value and
dominates all entries. Apply the checked column estimate.
\item \textbf{Grid membership.} If $-R_k\le z_k\le R_k$, the index
$z_k+R_k$ belongs to the coordinate range of length $2R_k+1$.
Prove that index returns $z_k/\Delta_k$. Cartesian membership follows
by finite list induction. This establishes \code{GridCompletenessGoal}.
\item \textbf{Filtering.} The checked theorem
\code{enumeration\_of\_gridCompleteness} already derives
\code{EnumerationGoal} from the previous target. Duplicate-freedom
is a separate short consequence of injectivity of each fixed-denominator
coordinate map, followed by product and filter preservation.
\end{enumerate}
The only premise of the final filter theorem is the location theorem,
not the full desired enumeration assertion under a new name. Every
mathematical step needed for that location theorem is proved in the
human argument above.

\subsection{From the host specification to object arithmetic}

The final embedded theorem uses numerical codes, not Lean's types
\code{Rat}, \code{Vector} or arbitrary predicates. Let
$\CertCheck(F,C)$, $Z(F,p)$ and $E(F,C)$ be the specified PR checker,
zero test and list-producing program. With validity and arity tests
included, the target is the first-order sentence
\[
 \forall F,C,p\;\bigl(\CertCheck(F,C)=1\ \land\ \operatorname{Valid}(p)
 \ \longrightarrow\
 [\operatorname{mem}(p,E(F,C))=1\ \leftrightarrow\ Z(F,p)=1]\bigr).
 \tag{E}\label{eq:object-enumeration}
\]
The syntax of the fixed programs is built from zero, successor,
projections, composition and primitive recursion. For each program,
prove its local equations in \code{Prf PRAAx}; then prove its loop
invariant with \code{PRA.qfInd}. The useful invariants are recorded
in Table~\ref{tab:invariants}.

Bounds are computed values of fixed PR functions; they need not be
small. A nested computation whose bound is obtained from an earlier
finite computation remains a PR composition. Proofs keep $F,C,p$
free and quantify only after the finite invariants are established.
This avoids an induction hypothesis containing an unbounded universal
claim about all rational zeros.

\begin{longtable}{P{.22\linewidth}P{.43\linewidth}P{.25\linewidth}}
\caption{Finite invariants needed for the full object derivation.}
\label{tab:invariants}\\
\toprule
Program & Quantifier-free content after bounded Booleanization & Fuel\\
\midrule\endfirsthead
\toprule Program & Invariant content & Fuel\\\midrule\endhead
Extended gcd & Current remainder equations and integer combinations
of the two original inputs. & Sum of input magnitudes plus one.\\[3pt]
Vector Bezout fold & The accumulated gcd is a linear combination of
the processed entries and divides them. & Vector length.\\[3pt]
Coefficient collection & Each processed monomial contributes its
coefficient to the matching exponent slot. & Input list length and
support size.\\[3pt]
Polynomial evaluation & The accumulator is the sum of the values of
the processed monomials. & List length, variable count and exponents.\\[3pt]
Matrix action & Each partial sum agrees with the processed column
terms at the fixed point. & Matrix dimension.\\[3pt]
Maximum scan & The selected index is processed, its value is the
current maximum, and all processed values are at most it. & Vector
length.\\[3pt]
Grid decoding & Each processed digit is in range and the quotient/
remainder reconstruction equation holds. & Number of coordinates.\\[3pt]
Filter & For each processed candidate, output membership is equivalent
to passing the exact equation tests. & Product of coordinate-grid sizes.\\
\bottomrule
\end{longtable}

Once \eqref{eq:object-enumeration} has a \code{Prf PRAAx}
derivation, the existing theorem \code{pra\_iff\_native} transports
it to the native scalar presentation \cite{NativeSource}. That bridge
is a final step. It cannot replace construction of the object derivation.

The printed axiom dependencies also distinguish the layers. The new
object builders and integer eigenvalue lemmas use only the host kernel
principles \code{propext} and \code{Quot.sound}; no
\code{Classical.choice} occurs there. The standard rational-number
implementation used for the host target programs does introduce
\code{Classical.choice} into the host filter theorem's dependencies.
That does not add a rule to \code{PRAAx}. The proposed arithmetic
implementation uses explicit signed-integer and positive-denominator
codes, so its constructive content is not inferred from that host
dependency list. Lean's acceptance of a recursion definition likewise
does not by itself prove primitive recursiveness or reify it into
\code{PR} syntax.
